\chapter{Оптимизация и основы машинного обучения}\label{chap:ml-foundations}\markright{Глава 49}
\section{Модель, параметры и измеримая цель}
Эта глава связывает линейную алгебру, производные и статистику. Все матрицы, входы и параметры далее действительные и конечные, размеры и число наблюдений --- положительные целые; для логарифмической потери используется натуральный логарифм.
\begin{definition}\term{Признак} (\emph{feature}) --- числовая характеристика объекта; список признаков образует вход \(x\in\mathbb R^d\).\end{definition}
\begin{definition}\term{Модель с параметрами} (\emph{parametric model}) --- выбранное семейство функций \(f_\theta\), где \(\theta\in\mathbb R^p\) --- изменяемые параметры.\end{definition}
\begin{definition}\term{Функция потерь} (\emph{loss function}) \(\ell(\widehat y,y)\) измеряет несоответствие предсказания \(\widehat y\) и наблюдаемого ответа \(y\).\end{definition}
\begin{definition}\term{Эмпирическая целевая функция} (\emph{empirical objective}) на \(N\ge1\) обучающих парах \((x_i,y_i)\) равна
\[J(\theta)=\frac1N\sum_{i=1}^N\ell(f_\theta(x_i),y_i).\]
\end{definition}
\par\addvspace{2pt}\begin{lrbox}{\CatFlowBox}\begin{minipage}[t]{\linewidth}
\setlength{\parskip}{2pt}\setlength{\columnsep}{4.5mm}\setlength{\intextsep}{0pt}
\begin{wrapfigure}[10]{r}{54mm}
\includegraphics[width=54mm,height=58mm,keepaspectratio]{images/unique/art-275.png}
\end{wrapfigure}
Обучение модели начинается с вопроса о том, какие величины известны и какую величину нужно предсказывать. Затем выбирают семейство формул и указывают изменяемые параметры. Данные сами по себе не выбирают ни семейство, ни способ измерять ошибку: это решения автора модели. Одна формула может хорошо описывать уже наблюдавшиеся точки и плохо работать на новых. Поэтому вычисление малого значения ошибки на знакомой таблице ещё не завершает задачу. Полезно отдельно записать допустимые входы, единицы измерения и ожидаемый вид ответа. После этого становится понятно, какие преобразования имеют смысл, какие параметры следует искать и по каким данным допустимо судить о результате. Математическая оптимизация отвечает на точно поставленный вопрос, но не исправляет неверную постановку.
\par\WFclear\end{minipage}\end{lrbox}%
\ifdim\dimexpr\ht\CatFlowBox+\dp\CatFlowBox\relax<58mm\dp\CatFlowBox=\dimexpr58mm-\ht\CatFlowBox\relax\fi%
\noindent\usebox{\CatFlowBox}\par\addvspace{2pt}
Для числового ответа удобно взять \(\ell(\widehat y,y)=\frac12(\widehat y-y)^2\). Остаток \(\widehat y-y\) может иметь знак, а эта потеря неотрицательна. Для меток \(y\in\{0,1\}\) применим другую потерю ниже. Выбор потери является частью постановки задачи, а не следствием названия алгоритма.
\section{Градиентный спуск: точный разбор квадратичной модели}
\begin{definition}\term{Выпуклое множество} (\emph{convex set}) \(D\subseteq\mathbb R^p\) вместе с любыми \(x,y\in D\) содержит весь отрезок между ними: \((1-t)x+ty\in D\) при \(0\le t\le1\).\end{definition}
\begin{definition}\term{Выпуклая функция} (\emph{convex function}) \(f:D\to\mathbb R\) на выпуклом множестве удовлетворяет
\[f((1-t)x+ty)\le(1-t)f(x)+tf(y)\qquad(x,y\in D,\ 0\le t\le1).\]
График на отрезке лежит не выше соединяющей его концы хорды.\end{definition}
\begin{definition}\term{Строго выпуклая функция} (\emph{strictly convex function}) удовлетворяет строгому неравенству в предыдущем определении при \(x\ne y\), \(0<t<1\).\end{definition}
\begin{example}Функция \(f(x)=x^2\) на \(\mathbb R\) строго выпукла, поскольку
\[(1-t)a^2+tb^2-((1-t)a+tb)^2=t(1-t)(a-b)^2>0\]
при \(a\ne b\), \(0<t<1\). Строгая выпуклость даёт не более одного минимизатора: если два разных минимума достигнуты, их середина имеет ещё меньшее значение. Существование минимума само по себе из строгой выпуклости не следует.\end{example}

\begin{definition}\term{Градиентный спуск} (\emph{gradient descent}) для дифференцируемой \(J\) выполняет шаг
\[\theta_{k+1}=\theta_k-\eta\nabla J(\theta_k),\qquad \eta>0.\]
Число \(\eta\) называют \term{шагом обучения} (\emph{learning rate}). Все компоненты одного шага вычисляют по прежнему вектору параметров.\end{definition}
\par\addvspace{2pt}\begin{lrbox}{\CatFlowBox}\begin{minipage}[t]{\linewidth}
\setlength{\parskip}{2pt}\setlength{\columnsep}{4.5mm}\setlength{\intextsep}{0pt}
\begin{wrapfigure}[12]{l}{54mm}
\includegraphics[width=54mm,height=58mm,keepaspectratio]{images/unique/art-276.png}
\end{wrapfigure}
Шаг оптимизации полезно рассматривать как проверяемое изменение параметров, а не как обещание улучшения при любом выборе числа. В квадратичной модели можно точно вычислить, во сколько раз изменится отклонение от решения. Слишком большой шаг способен перескочить через минимум и увеличить ошибку, хотя направление производной вычислено правильно. Для нескольких координат одно и то же число шага действует на направления с разной кривизной. Поэтому удобный масштаб признаков и разумный выбор шага связаны между собой. При вычислениях стоит сохранять прежние значения всех параметров до окончания одного шага, иначе получится другая последовательность действий. Простейший пример помогает увидеть эти различия до перехода к сложной модели, где рисунок поверхности ошибки уже трудно представить.
\par\WFclear\end{minipage}\end{lrbox}%
\ifdim\dimexpr\ht\CatFlowBox+\dp\CatFlowBox\relax<58mm\dp\CatFlowBox=\dimexpr58mm-\ht\CatFlowBox\relax\fi%
\noindent\usebox{\CatFlowBox}\par\addvspace{2pt}
\begin{theorem}[Шаг для скалярной квадратичной функции]\label{thm:review-ml-gd}
Пусть \(f(t)=\frac a2(t-b)^2\), \(a>0\). При постоянном \(0<\eta<\frac2a\) градиентный спуск из любого \(t_0\) сходится к единственному минимуму \(b\). Если \(t_k\ne b\), значение функции на каждом шаге строго уменьшается. При \(\eta=\frac2a\) ненулевое отклонение меняет знак без уменьшения модуля; при \(\eta>\frac2a\) его модуль растёт.
\end{theorem}
\begin{proof}Поскольку \(f^{\prime}(t)=a(t-b)\), получаем
\[t_{k+1}-b=(1-\eta a)(t_k-b),\qquad t_k-b=(1-\eta a)^k(t_0-b).\]
При \(|1-\eta a|<1\) геометрическая последовательность стремится к нулю, а \(f(t_{k+1})=(1-\eta a)^2f(t_k)\). Условие на модуль равносильно \(0<\eta<\frac2a\). При \(\eta=\frac2a\) множитель равен \(-1\), при большем шаге меньше \(-1\). Если начальная точка уже минимум, она остаётся им при любом шаге.\end{proof}
\begin{example}Для \(F(x,y)=\frac12x^2+2y^2\) градиент равен \((x,4y)\). Ошибки двух координат умножаются на \(1-\eta\) и \(1-4\eta\), поэтому общий безопасный диапазон для сходимости из любой точки --- \(0<\eta<\frac12\). При \(\eta=\frac14\) из \((2,1)\) получаем \((\frac32,0)\), а значение уменьшается с \(4\) до \(\frac98\). При \(\eta=1\) вторая координата умножается на \(-3\), поэтому при ненулевом начальном \(y\) метод расходится.\end{example}
Доказанное условие относится к указанной квадратичной модели. Для произвольной гладкой функции, тем более для нейросети, здесь не утверждается ни сходимость при любом малом постоянном шаге, ни нахождение глобального минимума. Малый градиент сам по себе также не доказывает, что точка является минимумом.
\section{Линейная модель как задача наименьших квадратов}
\begin{definition}\term{Линейная регрессия} (\emph{linear regression}) в простейшем виде использует \(f_{w,b}(x)=w^{\mathrm T}x+b\). Зависимость от \(x\) при \(b\ne0\) аффинна; слово «линейная» здесь относится к линейности по параметрам после добавления постоянного признака \(1\).\end{definition}
Добавим к каждому входу последнюю координату \(1\) и составим строки матрицы \(A\). При \(\theta=(w,b)\) и столбце ответов \(y\)
\[J(\theta)=\frac1{2N}\|A\theta-y\|^2,\qquad \nabla J(\theta)=\frac1N A^{\mathrm T}(A\theta-y).\]
Для проверки градиента раскроем \(J(\theta+h)-J(\theta)\): линейная по \(h\) часть равна \(\frac1N h^{\mathrm T}A^{\mathrm T}(A\theta-y)\), остаток --- \(\frac1{2N}\|Ah\|^2\). Минимизаторы удовлетворяют нормальным уравнениям из главы \ref{chap:c26}. Коэффициенты единственны при независимости столбцов, а при зависимых столбцах одни предсказания могут задаваться разными коэффициентами.
\begin{example}Для данных \((0,1),(1,2),(2,2)\) матрица имеет строки \((0,1),(1,1),(2,1)\). Нормальные уравнения \(5w+3b=6\), \(3w+3b=5\) дают \(w=\frac12,b=\frac76\). Ошибки предсказания \((\frac16,-\frac13,\frac16)\) имеют сумму квадратов \(\frac16\), поэтому \(J=\frac1{36}\). Множитель перед суммой квадратов не меняет точку минимума, но меняет градиент и допустимый масштаб шага.\end{example}
\section{Сигмоида и бинарная логистическая потеря}
\par\addvspace{2pt}\begin{lrbox}{\CatFlowBox}\begin{minipage}[t]{\linewidth}
\setlength{\parskip}{2pt}\setlength{\columnsep}{4.5mm}\setlength{\intextsep}{0pt}
\begin{wrapfigure}[13]{l}{54mm}
\includegraphics[width=54mm,height=58mm,keepaspectratio]{images/unique/art-293.png}
\end{wrapfigure}
Для ответа с двумя возможными метками модель может сначала вычислить произвольное действительное число, а затем перевести его в число между нулём и единицей. Так разделяются линейный расчёт и выбранная вероятностная модель. Значение до преобразования не является вероятностью. Порог для принятия решения также не заменяет функцию потерь: два предсказания по одну сторону порога могут иметь разные ошибки обучения. Логарифмическая потеря учитывает не только совпадение итоговой метки, но и уверенность выбранной модели. При выводе производной важно сохранять зависимость вероятности от исходного числа. Тогда цепное правило сокращает выражение гораздо сильнее, чем отдельное дифференцирование несвязанных частей.
\par\WFclear\end{minipage}\end{lrbox}%
\ifdim\dimexpr\ht\CatFlowBox+\dp\CatFlowBox\relax<61mm\dp\CatFlowBox=\dimexpr61mm-\ht\CatFlowBox\relax\fi%
\noindent\usebox{\CatFlowBox}\par\addvspace{2pt}
\relax

\begin{definition}\term{Сигмоида, или логистическая функция} (\emph{sigmoid; logistic function}) задаётся \(\sigma(z)=\frac1{1+e^{-z}}\), \(z\in\mathbb R\); её значения строго между \(0\) и \(1\).\end{definition}
\begin{definition}\term{Логит модели} (\emph{model logit}) --- число \(z=w^{\mathrm T}x+b\); логистическая модель выдаёт \(p=\sigma(z)\). Его можно интерпретировать как оценку вероятности метки \(1\), но качество такой оценки требует проверки по данным.\end{definition}
\begin{definition}\term{Бинарная логистическая потеря} (\emph{binary logistic loss; binary cross-entropy}) при \(y\in\{0,1\}\) равна
\[\ell(z,y)=-y\ln\sigma(z)-(1-y)\ln(1-\sigma(z)).\]
\end{definition}
При конечном \(z\) оба аргумента логарифма положительны. Для численной работы лучше не вычислять логарифм уже округлённого значения \(p\): оно может стать машинным нулём или единицей.
\begin{definition}\term{Мягкий плюс} (\emph{softplus}) --- функция \(s(t)=\ln(1+e^t)\). Устойчивая равносильная запись:
\[s(t)=\max(t,0)+\ln(1+e^{-|t|}),\qquad \ell(z,y)=s((1-2y)z).\]
\end{definition}
В этой записи показатель экспоненты неположителен, что предотвращает её переполнение; для малого положительного \(u\) функцию \(\ln(1+u)\) в программах вычисляют специальной операцией \(\operatorname{log1p}(u)\). Сигмоиду также можно вычислять по частям: \(\frac1{1+e^{-z}}\) при \(z\ge0\), \(\frac{e^z}{1+e^z}\) при \(z<0\). Это улучшает устойчивость вычисления, но не отменяет конечную точность и возможное округление очень малых чисел до нуля.
\begin{theorem}[Производная логистической потери]\label{thm:review-ml-logistic}
Для конечного \(z\in\mathbb R\), \(y\in\{0,1\}\)
\[\sigma^{\prime}(z)=\sigma(z)(1-\sigma(z)),\qquad \frac{\partial\ell}{\partial z}=\sigma(z)-y,\qquad \frac{\partial^2\ell}{\partial z^2}=\sigma(z)(1-\sigma(z))>0.\]
\end{theorem}
\begin{proof}Дифференцирование \((1+e^{-z})^{-1}\) даёт \(\frac{e^{-z}}{(1+e^{-z})^2}\), то есть первую формулу. Подстановка определения сигмоиды в потерю даёт \(\ell(z,y)=\ln(1+e^z)-yz\). Производная равна \(\frac{e^z}{1+e^z}-y=\sigma(z)-y\); следующая производная положительна. Равенство с устойчивой записью проверяется отдельно при \(y=0,1\) и \(t\ge0,t<0\). Здесь дифференцируется гладкая функция, хотя её вычислительная запись содержит модуль и максимум.\end{proof}
\begin{example}Если \(z=\ln3\), то \(p=\frac34\). При \(y=1\) потеря \(\ln\frac43\), производная по \(z\) равна \(-\frac14\); при \(y=0\) потеря \(\ln4\), производная \(\frac34\). Для одного признака \(\frac{\partial\ell}{\partial w}=(p-y)x\), \(\frac{\partial\ell}{\partial b}=p-y\). Эти две производные относятся к параметрам, а не к входным меткам.\end{example}
Даже выпуклая задача не обязана иметь конечный минимизатор: для единственного наблюдения \(y=1,z=\theta\) потеря \(\ln(1+e^{-\theta})\) стремится к нулю при \(\theta\to+\infty\), но строго положительна при каждом конечном \(\theta\). Поэтому строгая выпуклость по логиту не является обещанием существования конечных оптимальных параметров на любой таблице.
\section{Слои и обратное распространение производных}
\begin{definition}\term{Аффинный слой} (\emph{affine layer}) вычисляет \(z=Wx+b\), где \(W\in\mathbb R^{m\times d}\), \(x\in\mathbb R^d\), \(b\in\mathbb R^m\); выход имеет \(m\) координат.\end{definition}
\begin{definition}\term{Выпрямленная линейная функция} (\emph{rectified linear unit; ReLU}) равна \(\rho(t)=\max(0,t)\). К вектору её применяют по координатам. При \(t<0\) производная \(0\), при \(t>0\) --- \(1\), в нуле обычной производной нет.\end{definition}
\begin{definition}\term{Нейронная сеть прямого распространения} (\emph{feedforward neural network}) --- композиция аффинных слоёв и выбранных нелинейных функций. В простом случае \(f(x)=W_2\rho(W_1x+b_1)+b_2\).\end{definition}
Без нелинейности композиция аффинных слоёв снова аффинна. Нелинейная функция меняет семейство доступных зависимостей, но увеличение числа слоёв само по себе не гарантирует ни хорошего обучения, ни хорошего предсказания.
\begin{definition}\term{Обратное распространение производных} (\emph{backpropagation}) вычисляет производные итоговой потери по параметрам, передавая их по вычислительному графу в обратном порядке с помощью цепного правила. Это способ вычислить градиент, а не отдельное правило обновления параметров.\end{definition}
\par\addvspace{2pt}\begin{lrbox}{\CatFlowBox}\begin{minipage}[t]{\linewidth}
\setlength{\parskip}{2pt}\setlength{\columnsep}{4.5mm}\setlength{\intextsep}{0pt}
\begin{wrapfigure}[11]{r}{54mm}
\includegraphics[width=54mm,height=58mm,keepaspectratio]{images/unique/art-277.png}
\end{wrapfigure}
Обратный проход по вычислениям начинается с того, как итоговая ошибка зависит от последнего результата. Затем эта чувствительность передаётся к более ранним промежуточным величинам по цепному правилу. Полезно сначала подписать все значения прямого прохода: без них легко перепутать вход слоя с его выходом. Каждый множитель в производной должен иметь понятный источник, а при нескольких путях влияния вклады складываются. Для одного скалярного примера можно выписать все шаги вручную и проверить результат небольшим изменением параметра. Такая проверка помогает обнаружить ошибку знака или пропущенный множитель, но не заменяет формулу производной. В точках излома требуется отдельное соглашение вычислительной процедуры; обычное цепное правило там нельзя применять без оговорки о существовании производных.
\par\WFclear\end{minipage}\end{lrbox}%
\ifdim\dimexpr\ht\CatFlowBox+\dp\CatFlowBox\relax<58mm\dp\CatFlowBox=\dimexpr58mm-\ht\CatFlowBox\relax\fi%
\noindent\usebox{\CatFlowBox}\par\addvspace{2pt}
Рассмотрим скалярную сеть
\[z=wx+b,\quad h=\rho(z),\quad s=vh+c,\quad L=\frac12(s-y)^2.\]
Вне \(z=0\) положим \(\delta_s=s-y\), \(\delta_z=\delta_s v\rho^{\prime}(z)\). Цепное правило даёт
\[
\frac{\partial L}{\partial v}=\delta_s h,\quad
\frac{\partial L}{\partial c}=\delta_s,\quad
\frac{\partial L}{\partial w}=\delta_z x,\quad
\frac{\partial L}{\partial b}=\delta_z.
\]
\begin{example}Возьмём \(x=2,y=1,w=1,b=-1,v=3,c=0\). Прямой проход: \(z=1,h=1,s=3,L=2\). Обратный проход: \(\delta_s=2\), поэтому \(\partial_vL=2\), \(\partial_cL=2\). Поскольку \(z>0\), \(\rho^{\prime}(z)=1\), значит, \(\delta_z=2\cdot3=6\), \(\partial_wL=6\cdot2=12\), \(\partial_bL=6\). Например, цепочка для \(w\) имеет четыре множителя:
\[\frac{\partial L}{\partial w}=\frac{\partial L}{\partial s}\frac{\partial s}{\partial h}\frac{\partial h}{\partial z}\frac{\partial z}{\partial w}=2\cdot3\cdot1\cdot2=12.\]
Один одновременный шаг \(\eta=\frac1{100}\) даёт \(w=\frac{22}{25}\), \(b=-\frac{53}{50}\), \(v=\frac{149}{50}\), \(c=-\frac1{50}\). Новый прямой проход: \(z=\frac7{10}>0\), \(s=\frac{1033}{500}\), \(L=\frac{284089}{500000}<2\). Уменьшение подтверждено для этого шага; из него не следует сходимость всех следующих шагов.\end{example}
В точке \(z=0\) программы могут договориться использовать для ReLU число \(0\) при обратном проходе. Это выбранное правило, а не значение существующей обычной производной. При \(z<0\) производные по \(w,b\) в этой сети нулевые, но производная по \(c\) может оставаться ненулевой.
\section{Матричный обратный проход}
Во всех формулах этого раздела векторы записываются столбцами. Пусть \(d,m\) --- положительные целые числа,
\[
 x\in\mathbb R^d,\qquad W\in\mathbb R^{m\times d},\qquad
 b,z\in\mathbb R^m,\qquad z=Wx+b.
\]
Если \(L\) зависит от матрицы \(W\), то её \underline{матричный градиент} (matrix gradient) состоит из частных производных по элементам:
\[
 (\nabla_W L)_{ij}=\frac{\partial L}{\partial W_{ij}}.
\]
Такой градиент имеет тот же размер \(m\times d\), что и \(W\).
Для приращения \(H\) матрицы \(W\) линейная часть приращения потери равна
\(\sum_{i=1}^m\sum_{j=1}^d(\nabla_WL)_{ij}H_{ij}\).
Это обычный полный дифференциал по \(md\) вещественным координатам; нового понятия производной здесь не требуется.

\begin{theorem}[Градиенты аффинного слоя]
Пусть функция \(\ell:\mathbb R^m\to\mathbb R\) дифференцируема в точке \(z=Wx+b\),
\(L=\ell(Wx+b)\), и
\[
 g=\nabla_z\ell(z)\in\mathbb R^m.
\]
Тогда, при вычислении каждой производной с остальными аргументами фиксированными,
\[
 \nabla_x L=W^{\mathrm T}g\in\mathbb R^d,\qquad
 \nabla_W L=gx^{\mathrm T}\in\mathbb R^{m\times d},\qquad
 \nabla_b L=g\in\mathbb R^m.
\]
\end{theorem}
\begin{proof}
При изменении только \(x\) на приращение \(u\) линейная часть приращения \(L\)
равна \(g^{\mathrm T}Wu=(W^{\mathrm T}g)^{\mathrm T}u\).
При изменении только \(b\) она равна \(g^{\mathrm T}v\), где \(v\) --- приращение \(b\).
При изменении \(W\) на матрицу \(H\) линейная часть равна
\[
 g^{\mathrm T}Hx
 =\sum_{i=1}^m\sum_{j=1}^d g_i x_j H_{ij}.
\]
Значит, \(\frac{\partial L}{\partial W_{ij}}=g_i x_j\), что и даёт матрицу
\(gx^{\mathrm T}\). Все три формулы следуют из полного правила цепочки
главы~\ref{chap:multivariable}.
\end{proof}

Теперь пусть \(h=\rho(z)\), где ReLU применяется к каждой координате,
а \(L=\ell(h)\). Для двух векторов одинакового размера запись
\(u\odot v\) означает их \underline{покоординатное произведение} (elementwise product):
\((u\odot v)_i=u_iv_i\). Если \(z_i\ne0\) для всех \(i\), положим
\[
 r=\nabla_h\ell(h)\in\mathbb R^m,\qquad
 a_i=\begin{cases}1,&z_i>0,\\0,&z_i<0,\end{cases}\qquad
 \delta=r\odot a\in\mathbb R^m.
\]
Тогда \(\delta=\nabla_z L\), и в формулах теоремы нужно заменить \(g\) на \(\delta\).
В частности,
\[
 \nabla_x L=W^{\mathrm T}\delta,\qquad
 \nabla_W L=\delta x^{\mathrm T},\qquad
 \nabla_b L=\delta.
\]
В нуле у самой функции ReLU обычной производной нет. Если какой-то \(z_i=0\),
эти формулы с выбранной программной маской уже не являются применением
классического правила цепочки к дифференцируемой ReLU в этой координате.
У всей композиции производная в особом случае может существовать, но это
нужно проверять отдельно.

Для двух слоёв
\[
 z=W_1x+b_1,\qquad h=\rho(z),\qquad s=W_2h+b_2,\qquad L=\ell(s),
\]
размеры равны
\(W_1\in\mathbb R^{m\times d}\), \(W_2\in\mathbb R^{k\times m}\),
\(b_1\in\mathbb R^m\), \(b_2,s\in\mathbb R^k\).
При \(z_i\ne0\) обратный проход имеет вид
\[
 \delta_2=\nabla_s\ell(s)\in\mathbb R^k,\qquad
 \delta_1=(W_2^{\mathrm T}\delta_2)\odot a\in\mathbb R^m.
\]
Отсюда
\[
 \nabla_{W_2}L=\delta_2h^{\mathrm T},\quad \nabla_{b_2}L=\delta_2,\qquad
 \nabla_{W_1}L=\delta_1x^{\mathrm T},\quad \nabla_{b_1}L=\delta_1,\qquad
 \nabla_xL=W_1^{\mathrm T}\delta_1.
\]
Это два последовательных применения доказанной теоремы и правила цепочки.
Все величины обратного прохода вычисляются по прежним параметрам.
Обновление параметров градиентным шагом выполняется после этого.

\begin{example}
Возьмём
\[
 x=\begin{pmatrix}1\\-2\end{pmatrix},\qquad
 W=\begin{pmatrix}1&1\\-1&1\\2&0\end{pmatrix},\qquad
 b=\begin{pmatrix}2\\0\\-1\end{pmatrix},\qquad
 y=\begin{pmatrix}0\\1\\-1\end{pmatrix}.
\]
Пусть \(h=\rho(Wx+b)\), а \(L=\frac{1}{2}\|h-y\|^2\).
Прямой проход даёт
\[
 z=\begin{pmatrix}1\\-3\\1\end{pmatrix},\qquad
 h=\begin{pmatrix}1\\0\\1\end{pmatrix},\qquad L=3.
\]
Градиент по \(h\) равен \(r=h-y=(1,-1,2)^{\mathrm T}\), маска
\(a=(1,0,1)^{\mathrm T}\), поэтому \(\delta=(1,0,2)^{\mathrm T}\).
Обратный проход:
\[
 \nabla_xL=\begin{pmatrix}5\\1\end{pmatrix},\qquad
 \nabla_WL=\begin{pmatrix}1&-2\\0&0\\2&-4\end{pmatrix},\qquad
 \nabla_bL=\begin{pmatrix}1\\0\\2\end{pmatrix}.
\]
Умножение \(W^{\mathrm T}\delta\) согласовано по размерам:
матрица \(2\times3\) переводит столбец из \(\mathbb R^3\) в столбец из \(\mathbb R^2\).
Средняя строка матричного градиента нулевая, потому что вторая координата ReLU не активна.
\end{example}
\par\addvspace{2pt}\begin{lrbox}{\CatFlowBox}\begin{minipage}[t]{\linewidth}
\setlength{\parskip}{2pt}\setlength{\columnsep}{4.5mm}\setlength{\intextsep}{0pt}
\begin{wrapfigure}[12]{r}{54mm}
\includegraphics[width=54mm,height=58mm,keepaspectratio]{images/unique/art-294.png}
\end{wrapfigure}
В матричном обратном проходе размерности помогают проверить каждое умножение ещё до подстановки чисел. Градиент по вектору записываем столбцом той же длины, что и этот вектор; градиент по матрице имеет столько же строк и столбцов, сколько параметр. Прямое вычисление слоя переводит вход в пространство выходов, а обратное передаёт влияние потери в обратном направлении через транспонированную матрицу. Это не обращение исходного отображения: обратимая матрица здесь не требуется. Если один параметр влияет на несколько выходов, соответствующие вклады складываются. Покоординатная нелинейность добавляет множители для каждой координаты, причём в точках излома обычное дифференцирование требует отдельной оговорки.
\par\WFclear\end{minipage}\end{lrbox}%
\ifdim\dimexpr\ht\CatFlowBox+\dp\CatFlowBox\relax<58mm\dp\CatFlowBox=\dimexpr58mm-\ht\CatFlowBox\relax\fi%
\noindent\usebox{\CatFlowBox}\par\addvspace{2pt}
\relax
\section{Обучение, проверка и масштаб признаков}
\begin{definition}\term{Обучающая выборка} (\emph{training set}) используется для подбора параметров.\end{definition}
\begin{definition}\term{Валидационная выборка} (\emph{validation set}) используется для выбора настроек процедуры: например, шага, числа шагов или размера модели. Такие настройки называют \term{гиперпараметрами} (\emph{hyperparameters}).\end{definition}
\begin{definition}\term{Тестовая выборка} (\emph{test set}) сохраняется для заключительной оценки выбранной процедуры; систематически подбирать по ней модель означает использовать её для настройки.\end{definition}
\begin{definition}\term{Переобучение} (\emph{overfitting}) --- ситуация, когда подгонка особенностей обучающих данных не даёт сопоставимого качества на новых данных из интересующей задачи.\end{definition}
Малая обучающая ошибка сама по себе не исключает переобучения. Валидация помогает сравнивать варианты, но многократный выбор по одной и той же таблице тоже может подгонять её особенности. Разбиение должно соответствовать задаче: связанные измерения одного объекта и данные из будущего нельзя бездумно перемешивать с прошлым. Такая утечка информации делает оценку качества слишком благоприятной.
\begin{definition}\term{Масштабирование признака} (\emph{feature scaling}) заменяет \(x\) на \(\frac{x-\mu}{s}\), где \(s>0\); часто \(\mu\) и \(s\) оценивают по обучающим данным.\end{definition}
Например, по обучающим значениям \(1,3\) среднее равно \(2\), корень из среднего квадрата отклонений равен \(1\); преобразованные значения \(-1,1\), а новый вход \(5\) преобразуется в \(3\). Параметры преобразования затем сохраняют для валидации и теста. При \(s=0\) признак постоянен на обучении: деление невозможно, нужно явно выбрать обработку, например исключить такой столбец. Масштабирование может облегчить оптимизацию, но меняет геометрию PCA и не гарантирует улучшения предсказаний. Перед заключением сравнивают одну и ту же заранее выбранную меру качества на подходящих проверочных данных.

\Needspace{7\baselineskip}\section*{Задачи}
\begingroup\setlength{\parskip}{2pt}
\noindent\begin{minipage}[t]{\linewidth}\textbf{49.1.} Для модели \(f_{w,b}(x)=wx+b\), \((w,b)=(2,1)\), данных \((0,1),(1,2)\) найди остатки и \(J=\dfrac1{2N}\sum_i(f(x_i)-y_i)^2\).\par\end{minipage}\par
\noindent\begin{minipage}[t]{\linewidth}\textbf{49.2.} Сделай один шаг для \(f(t)=2(t-3)^2\) из \(t_0=0\) с \(\eta=\dfrac14\). Сравни значения функции.\par\end{minipage}\par
\noindent\begin{minipage}[t]{\linewidth}\textbf{49.3.}\starmark{} Для \(f(t)=2t^2\) найди все постоянные положительные шаги, обеспечивающие сходимость из любой начальной точки. Что происходит при граничном шаге и \(t_0=1\)?\par\end{minipage}\par
\noindent\begin{minipage}[t]{\linewidth}\textbf{49.4.} Для \(F(x,y)=\dfrac12x^2+2y^2\) сделай шаг \(\eta=\dfrac14\) из \((2,1)\), вычисли новое значение и общий диапазон постоянных шагов сходимости из любой точки.\par\end{minipage}\par
\noindent\begin{minipage}[t]{\linewidth}\textbf{49.5.} Для той же функции сделай шаг \(\eta=\dfrac34\) из \((0,1)\). Объясни, почему верный градиент не гарантирует уменьшения при таком шаге.\par\end{minipage}\par
\par\addvspace{2pt}\begin{lrbox}{\CatFlowBox}\begin{minipage}[t]{\linewidth}
\setlength{\parskip}{2pt}\setlength{\columnsep}{4.5mm}\setlength{\intextsep}{0pt}
\begin{wrapfigure}[12]{r}{54mm}
\includegraphics[width=54mm,height=58mm,keepaspectratio]{images/unique/art-278.png}
\end{wrapfigure}
\textbf{49.6.} Пусть \(A=\begin{pmatrix}1&1\\2&1\end{pmatrix}\), \(y=(1,2)^{\mathrm T}\), \(J=\dfrac14\|A\theta-y\|^2\). Найди \(J(0)\), градиент в нуле и шаг при \(\eta=\dfrac1{10}\).\par

\textbf{49.7.} Найди МНК-модель \(y=wx+b\) для двух точек \((0,1),(1,3)\). Почему коэффициенты единственны?\par

\textbf{49.8.}\starmark{} Для двух наблюдений \((1,1),(1,3)\) найди все МНК-модели \(wx+b\) и минимум \(J=\dfrac14\sum_i(wx_i+b-y_i)^2\).\par

\textbf{49.9.} Вычисли \(\sigma(0),\sigma(\ln3),\sigma(-\ln3)\).\par
\par\WFclear\end{minipage}\end{lrbox}%
\ifdim\dimexpr\ht\CatFlowBox+\dp\CatFlowBox\relax<60mm\dp\CatFlowBox=\dimexpr60mm-\ht\CatFlowBox\relax\fi%
\noindent\usebox{\CatFlowBox}\par\addvspace{2pt}

\noindent\begin{minipage}[t]{\linewidth}\textbf{49.10.} При \(z=\ln3\) найди \(\sigma^{\prime}(z)\) и производные логистической потери для меток \(1\) и \(0\).\par\end{minipage}\par
\noindent\begin{minipage}[t]{\linewidth}\textbf{49.11.} Вычисли точно \(\ell(0,1)\), \(\ell(\ln3,1)\), \(\ell(\ln3,0)\).\par\end{minipage}\par
\noindent\begin{minipage}[t]{\linewidth}\textbf{49.12.}\starmark{} Запиши устойчиво потери при \(z=1000\) для обеих бинарных меток. Почему вычисление через сначала округлённую сигмоиду может дать неверный бесконечный результат?\par\end{minipage}\par
\noindent\begin{minipage}[t]{\linewidth}\textbf{49.13.} Один обучающий объект имеет \(x=2,y=1\), модель \(z=wx+b\). При \(w=b=0\) найди градиент логистической потери и шаг \(\eta=\dfrac15\).\par\end{minipage}\par
\noindent\begin{minipage}[t]{\linewidth}\textbf{49.14.} Вычисли ReLU в точках \(-2,0,3\). Где из этих точек существует обычная производная и чему она равна?\par\end{minipage}\par
\noindent\begin{minipage}[t]{\linewidth}\textbf{49.15.}\starmark{} Докажи, что \(J(\theta)=\ln(1+e^{-\theta})\) на \(\mathbb R\) имеет инфимум \(0\), но не имеет минимизатора. Какой вывод о существовании обученных параметров нельзя делать из одной выпуклости?\par\end{minipage}\par
\noindent\begin{minipage}[t]{\linewidth}\textbf{49.16.} В сети \(z=wx+b,h=\rho(z),s=vh+c,L=\dfrac12(s-y)^2\) даны \(x=-1,y=0,w=2,b=3,v=-2,c=1\). Выполни прямой проход.\par\end{minipage}\par
\noindent\begin{minipage}[t]{\linewidth}\textbf{49.17.}\starmark{} Для той же архитектуры возьми \(x=1,y=0,w=2,b=1,v=-1,c=2\). Найди все четыре производные потери по параметрам.\par\end{minipage}\par
\noindent\begin{minipage}[t]{\linewidth}\textbf{49.18.} Теперь \(x=1,y=2,w=-1,b=0,v=3,c=1\). Найди прямой проход и четыре производные. Все ли параметры обязательно перестают изменяться?\par\end{minipage}\par
\noindent\begin{minipage}[t]{\linewidth}\textbf{49.19.}\starmark{} Докажи, что два согласованных по размерам аффинных слоя без нелинейности образуют один аффинный слой. Проверь на \(u=2x+1\), \(s=3u-4\).\par\end{minipage}\par
\noindent\begin{minipage}[t]{\linewidth}\textbf{49.20.} По обучающим значениям признака \(2,4\) используют \(\mu\) --- среднее и \(s\) --- корень из среднего квадрата отклонений. Преобразуй новый вход \(10\). Варианты A и B имеют ошибки обучения \(0{,}01,0{,}05\), валидации \(0{,}5,0{,}1\). Какой вариант предпочтительнее по указанной валидационной мере? Можно ли перенастроить масштаб по тестовым данным?\par\end{minipage}\par
\noindent\begin{minipage}[t]{\linewidth}\textbf{49.21.} Пусть \(x=(2,-1,3)^{\mathrm T}\), \(W=\begin{pmatrix}1&2&0\\-1&0&1\end{pmatrix}\), \(b=(1,-2)^{\mathrm T}\), \(y=(0,1)^{\mathrm T}\), \(z=Wx+b\), \(L=\dfrac{1}{2}\|z-y\|^2\). Найди прямой проход, потерю и градиенты по \(x,W,b\). Укажи их размеры.\par\end{minipage}\par
\noindent\begin{minipage}[t]{\linewidth}\textbf{49.22.}\starmark{} В двухслойной сети \(z=W_1x+b_1\), \(h=\rho(z)\), \(s=W_2h+b_2\), \(L=\dfrac{1}{2}(s-y)^2\) даны
\[
 x=(1,-1)^{\mathrm T},\quad W_1=\begin{pmatrix}1&0\\0&1\end{pmatrix},\quad
 b_1=(0,0)^{\mathrm T},\quad W_2=\begin{pmatrix}2&-1\end{pmatrix},\quad b_2=1,\quad y=1.
\]
Выполни прямой и обратный проходы. Найди градиенты по четырём параметрам \(W_1,b_1,W_2,b_2\) и по входу \(x\).\par\end{minipage}\par
\endgroup
